% Movimento Dependente: 2 blocos e 2 roldanas - auxiliar
% Faro, dom 07 dez 2025 10:16:53 WET
% Written by Tordar 
% pstricks2pdf.sh 2blocos_2roldanas02
\documentclass{article}
\pagestyle{empty}

\usepackage{pst-all} 
\usepackage{pst-slpe}
\usepackage{graphicx}
\input{apnum}

% Roldana
\newcommand{\roldana}[5]{% x, y, radius, hole-radius, color
 \pscircle[linewidth=0.5mm,linecolor=lightgray,dimen=middle,fillstyle=radslope,slopebegin=gray,slopeend=white,slopesteps=50](#1,#2){#3}
  \ifnum\pdfstrcmp{#4}{0}=0\relax\else
    \pscircle[fillcolor=white,fillstyle=solid,linewidth=0.8pt](#1,#2){#4}
  \fi
  \pscircle[linewidth=1pt,linecolor=#5!50!black](#1,#2){#3}
}
% Bloco
\newcommand{\bloco}[5][]{% #1=options, #2=cx, #3=cy, #4=totalWidth, #5=height
  % Left frame (half of total width, full height)
  \psframe[linewidth=1pt,linecolor=gray,fillstyle=slope,slopebegin=white,slopeend=gray,slopesteps=15]
    (!#2 #4 2 div sub #3 #5 2 div sub)                     % lower-left of total area
    (!#2 #3 #5 2 div add)                                  % upper-right of left frame (center x, top)  
  % Right frame (half of total width, full height)
  \psframe[linewidth=1pt,linecolor=gray,fillstyle=slope,slopebegin=gray,slopeend=white,slopesteps=15]
    (!#2 #3 #5 2 div sub)                                  % lower-left of right frame (center x, bottom)
    (!#2 #4 2 div add #3 #5 2 div add)                     % upper-right of total area
  \psframe[#1](!#2 #4 2 div sub #3 #5 2 div sub)(!#2 #4 2 div add #3 #5 2 div add) 
}

\begin{document}

% Raio de cada roldana:
\evaldef\ther{0.75}
% Altura e largura de cada bloco: 
\evaldef\thew{1.5}
\evaldef\theh{1.5}
% Coordenadas do primeiro bloco: 
\evaldef\xbi{3}
\evaldef\ybi{3}
% Posição da primeira roldana: 
\evaldef\xri{\xbi+\ther}
\evaldef\yri{7}
% Posição da segunda roldana: 
\evaldef\xrii{\xbi+3*\ther}
\evaldef\yrii{5}
% Coordenadas do segundo bloco: 
\evaldef\xbii{\xrii}
\evaldef\ybii{1}
% Coordenadas auxiliares:
\evaldef\yfi{\ybi+\theh/2}
\evaldef\xfii{\xrii-\ther}
\evaldef\yfii{\yrii}
\evaldef\xfiii{\xrii+\ther}
\evaldef\yfiv{\ybii+\theh/2}
\evaldef\xfv{\xrii+\ther}
\evaldef\xa{\xbi-0.75*\thew}
\evaldef\xb{\xbii+0.75*\thew}
\evaldef\xaf{\xbi-0.2}
\evaldef\yaf{\ybi+0.6*\theh}
\evaldef\xbf{\xaf}
\evaldef\ybf{\yri}
\evaldef\xcf{\xri+\ther+0.2}
\evaldef\ycf{\yri}
\evaldef\xdf{\xrii-\ther-0.2}
\evaldef\ydf{\yrii}
\evaldef\xef{\xrii+\ther+0.2}
\evaldef\yef{\yrii}
\evaldef\xff{\xef}
\evaldef\yff{8.7}
% Seta dupla 1:
\evaldef\xarri{\xri-\ther-0.8}
\evaldef\xxi{\xarri-0.4}
\evaldef\yxi{0.5*(\yri+9)}
% Seta dupla 2:
\evaldef\xarrii{\xrii+\ther+0.8}
\evaldef\yarrii{\ybii+\theh/2}
\evaldef\xxii{\xarrii+0.4}
\evaldef\yxii{0.5*(\yrii+\yarrii)}

% Let's go:
\begin{pspicture}(0,0)(10,10)
% Teto: 
\psframe[linewidth=0pt,linecolor=lightgray,dimen=inner,fillstyle=hlines*,fillcolor=lightgray,hatchangle=45](1.5,9)(8.5,10) 
% Blocos: 
\bloco[linewidth=2pt,linecolor=black]{\xbi}{\ybi}{\thew}{\theh} % 1
\bloco[linewidth=2pt,linecolor=black]{\xbii}{\ybii}{\thew}{\theh} % 2
% Roldanas e fios:
\roldana{\xri}{\yri}{\ther}{0.05}{gray} % 1 
\psline[linewidth=2pt,linecolor=black](\xri,\yri)(\xri,9)         % Fio teto      - roldana 1
\psline[linewidth=2pt,linecolor=black](\xbi,\yfi)(\xbi,\yri)      % Fio bloco   1 - roldana 1
\psline[linewidth=2pt,linecolor=black](\xfii,\yri)(\xfii,\yrii)   % Fio roldana 1 - roldana 2
\psline[linewidth=2pt,linecolor=black](\xfv,\yrii)(\xfv,9)        % Fio teto      - roldana 2
\roldana{\xrii}{\yrii}{\ther}{0.05}{gray} % 2
\psline[linewidth=2pt,linecolor=black](\xbii,\yfiv)(\xbii,\yrii) % Fio bloco 2   - roldana 2
% Texto: 
\rput{0}(\xa,\ybi){\scalebox{1.5}{$A$}}
\rput{0}(\xb,\ybii){\scalebox{1.5}{$B$}}
\rput{0}(\xaf,\yaf){\scalebox{1.2}{$a$}}
\rput{0}(\xbf,\ybf){\scalebox{1.2}{$b$}}
\rput{0}(\xcf,\ycf){\scalebox{1.2}{$c$}}
\rput{0}(\xdf,\ydf){\scalebox{1.2}{$d$}}
\rput{0}(\xef,\yef){\scalebox{1.2}{$e$}}
\rput{0}(\xff,\yff){\scalebox{1.2}{$f$}}
\rput{0}(\xxi,\yxi){\scalebox{1.5}{$x_1$}}
\rput{0}(\xxii,\yxii){\scalebox{1.5}{$x_2$}}
% Setas duplas: 
\psline[linewidth=2pt,linecolor=black]{<->}(\xarri,\yri)(\xarri,9) 
\psline[linewidth=2pt,linecolor=black]{<->}(\xarrii,\yarrii)(\xarrii,\yrii) 
\end{pspicture}

\end{document}
